\documentclass[10pt,a4paper]{article}
\usepackage{xeCJK}
\usepackage{fontspec}
\usepackage{booktabs}
\usepackage{array}
\usepackage{xcolor}
\usepackage{geometry}
\usepackage[protrusion=true]{microtype}
\geometry{margin=1.5cm}
\linespread{1.02}

\setmainfont{Baskerville}
\setCJKmainfont{Songti SC}
\newfontfamily{\starfont}{Songti SC}

\definecolor{wrongred}{RGB}{198,40,40}
\definecolor{rulegray}{RGB}{190,190,190}

\newcommand{\checkmarkbox}{\fbox{\rule{0pt}{0.75ex}\rule{0.75ex}{0pt}}}
\newcommand{\STAR}{\textcolor{wrongred}{\starfont ★}}
% \fa{题号}{答案}{解析}
\newcommand{\fa}[3]{\noindent\textbf{#1.}~\textbf{#2}\par\nobreak\vspace{0.02em}\hspace*{1.5em}{{\footnotesize #3}\par\vspace{0.12em}}}

\begin{document}
\pagestyle{plain}

\begin{center}
{\LARGE\bfseries 名词性从句\quad 参考答案与解析}\\[0.25em]
{\large 上海市大同中学高一英语 · 名词性从句专项（共 40 题）}\\[0.22em]
{\small 2029届高一英语 · 解析以“从句缺什么成分”为主线}
\end{center}

\vspace{0.3em}
\noindent{\small 说明：原卷未印答案，本答案按标准语法补拟；红笔判错记录不等于语法上必错，第 12、16、19、26、27、30、40 题的复核口径及可接受形式见各条解析。}

\vspace{0.4em}
\fa{1}{That}{主语从句本身完整，用 that 引导，不作成分、只起连接作用。}
\fa{2}{why}{That's why...（那就是……的原因）为固定句型；表结果用 why。}
\fa{3}{which}{有确定范围（so many options）中的“哪一个”，用 which。}
\fa{4}{what}{no longer 后为表语从句，表“它过去的样子”，what 作 used to be 的表语。}
\fa{5}{that}{suggestion 后为同位语从句，从句完整，用 that。}
\fa{6}{whether}{与 or not 连用只能用 whether，不能用 if。}
\fa{7}{What; that}{第一空 what 作 insisted on 的宾语并引导主语从句；第二空为表语从句，从句完整，用 that。}
\fa{8}{that}{be sure that...：宾语从句完整，用 that。}
\fa{9}{that}{it 为形式宾语，clear 后为真正的宾语从句，从句完整，用 that。}
\fa{10}{where}{表语从句表地点，用 where。}
\fa{11}{what}{It's none of your business what...：what 作 think about 的宾语（此句不能用 whether）。}
\fa{12}{why}{figure out 后为宾语从句，从句成分完整（insects, birds, and animals are dying out），缺的是原因状语 $\to$ why；学生答 if 被判错，说明本题不考“是否”。}
\fa{13}{where}{表语从句表地点（去哪儿找工作），用 where；题干无先行词可指代，不选 what。}
\fa{14}{who}{talk to 后缺宾语，指人用 who；how 表方式，与 to 的宾语位置冲突。}
\fa{15}{whoever}{介词 to 后为宾语从句，从句缺主语且表“无论谁”，用 whoever；不可用 if。}
\fa{16}{because}{表推测的原因：那是因为我在和校长谈话。\emph{when}（那时候）亦通，红笔按 because 口径判。}
\fa{17}{that}{no proof that...：同位语从句，从句完整，用 that。}
\fa{18}{what}{recognize 后缺宾语，且修饰 mistakes，用 what（what silly mistakes＝多么傻的错误）。}
\fa{19}{how}{表语从句表方式：如何说服他讲真话。\emph{whether}（是否）亦可，本题按红笔口径取 how。}
\fa{20}{Whoever}{主语从句缺主语，表“无论谁”，用 whoever；Who 只表疑问。}
\fa{21}{Who}{has won... hasn't been announced：疑问词引导主语从句表“谁”，用 who。}
\fa{22}{What; what}{第一空 what 作 get 的宾语；第二空比较对象为名词性从句，仍用 what，不可用 that。}
\fa{23}{That}{主语从句完整，用 that；句首 that 不可省。}
\fa{24}{which}{either... or... 已限定两者，用 which（哪一个）。}
\fa{25}{whatever}{take 后缺宾语、无先行词、表“任何你最喜欢的”，用 whatever。}
\fa{26}{Whether}{主语从句表“能否引起听众注意”，用 whether；raise 不与 how much 搭配。}
\fa{27}{whose}{表语从句中 name 前缺定语，用 whose（谁的名字）。}
\fa{28}{whatever}{give a child whatever he wants：从句缺宾语、表“任何他想要的”，用 whatever。}
\fa{29}{How}{主语从句表方式（你怎样学习工作），用 how。}
\fa{30}{whether}{thinks of 后表“是否能恰当应对”。本题存疑：\emph{how}（如何应对）语法亦完全成立，红笔却判 how 错（标记已放大复核，确为叉），故按红笔口径取 whether；两说并存，见订正表旁注（六）。}
\fa{31}{what}{look like 后缺宾语，用 what（what our village looks like）。}
\fa{32}{that}{A story goes that...：同位语从句，从句完整，用 that。}
\fa{33}{what}{at what I thought was a dangerous speed：what 作 was 的表语，I thought 为插入语。}
\fa{34}{whether}{同位语从句表“是否值得”，用 whether。}
\fa{35}{where}{表地点，用 where（最需要他的地方）。}
\fa{36}{what}{it takes 缺宾语，用 what（那就是做好任何事所需要的东西）。}
\fa{37}{what}{be surprised at 缺宾语，用 what 引导主语从句。}
\fa{38}{whatever}{do whatever lies in its power：缺宾语、表“任何力所能及的”，用 whatever。}
\fa{39}{that}{promise 后为同位语从句，从句完整，用 that。}
\fa{40}{Whoever's}{主语从句中 room 前缺定语，用 whoever's（＝anyone's who）。本题存疑：\emph{Whose room this is} 语法上亦成立，红笔判 Whose 错。}

\vspace{0.25em}
\noindent{\color{rulegray}\rule{\textwidth}{0.5pt}}
\vspace{0.25em}
\noindent\textbf{复核记录}\hfill\small 完成重做后填写

\vspace{0.25em}
\noindent\begin{tabular}{@{}p{0.31\textwidth}@{}p{0.34\textwidth}@{}p{0.35\textwidth}@{}}
重做得分：\underline{\hspace{1.2cm}}\,/\,40
& 仍错误的题号：\underline{\hspace{3.2cm}}
& 下次复习日期：\underline{\hspace{2.0cm}} \\
\end{tabular}

\vspace{0.25em}
\noindent 本轮易错点：\checkmarkbox\ 从句缺成分判断\quad\checkmarkbox\ what／whatever\quad\checkmarkbox\ whether／if\quad\checkmarkbox\ 连接代词 which／whoever

\vspace{0.2em}
\noindent 复习状态：\checkmarkbox\ 已掌握\quad\checkmarkbox\ 需要巩固\quad\checkmarkbox\ 需重点复习

\end{document}
